\documentclass[11pt,a4paper]{article}

\usepackage[utf8]{inputenc}
\usepackage[T1]{fontenc}
\usepackage{lmodern}
\usepackage{amsmath,amssymb}
\usepackage{booktabs}
\usepackage{graphicx}
\usepackage{hyperref}
\usepackage{geometry}
\usepackage{xcolor}

\geometry{margin=2.5cm}
\hypersetup{colorlinks=true,linkcolor=blue!50!black,citecolor=blue!50!black,urlcolor=blue!50!black}

\newcommand{\useritem}{(u,i)}
\newcommand{\Rset}{\mathcal{S}}
\newcommand{\Uset}{\mathcal{U}}
\newcommand{\Iset}{\mathcal{I}}
\newcommand{\pvec}{\mathbf{p}}
\newcommand{\qvec}{\mathbf{q}}

\title{\textbf{Bayesian Personalized Ranking for Anime Recommendation}\\[0.4em]
\large A reproducible implicit-feedback pipeline with point-in-time features\\
and an honest comparison against tuned neighbourhood baselines}

\author{Patrick Oliveira}
\date{\today}

\begin{document}

\maketitle

\begin{abstract}
We implement Bayesian Personalized Ranking (BPR) matrix factorization for
top-$N$ anime recommendation from implicit feedback, and benchmark it against
most-popular, item-$k$NN and random baselines under a leave-last-one-out
protocol. The contribution of this report is methodological rather than
algorithmic: every aggregate feature is computed \emph{point-in-time}, all
hyperparameters are selected on a validation partition that is disjoint from
the test partition, and the served ranking is asserted to be identical to the
ranking that was scored offline. Our headline result is negative for the
flagship model: a properly tuned item-$k$NN attains NDCG@10 of $0.1305$ versus
$0.1248$ for tuned BPR. We report this rather than tuning until the preferred
method wins, and we analyse the data regime that explains it.
\end{abstract}

\tableofcontents
\newpage

\section{Notation}
\label{sec:notation}

Let $\Uset$ be the set of users and $\Iset$ the set of items, with
$|\Uset| = m$ and $|\Iset| = n$. Implicit feedback is the set of observed
positive interactions
\begin{equation}
  \Rset \subseteq \Uset \times \Iset ,
\end{equation}
and $\Iset_u^{+} = \{\, i \in \Iset : (u,i) \in \Rset \,\}$ denotes the items
consumed by user $u$. A model produces a score $\hat{x}_{ui} \in \mathbb{R}$
used only for ranking; its absolute value carries no meaning.

Latent factors are written $\pvec_u \in \mathbb{R}^{d}$ for users and
$\qvec_i \in \mathbb{R}^{d}$ for items, collected in matrices
$P \in \mathbb{R}^{m \times d}$ and $Q \in \mathbb{R}^{n \times d}$. The full
parameter set is $\Theta = \{P, Q, \mathbf{b}\}$ with item biases
$\mathbf{b} \in \mathbb{R}^{n}$. The logistic function is
$\sigma(z) = (1 + e^{-z})^{-1}$.

\section{Introduction}
\label{sec:introduction}

Recommendation from implicit feedback differs from rating prediction in a way
that is easy to state and easy to get wrong. We never observe a negative
signal. A title a user has not watched may be a title they would dislike, or
simply one they have not encountered. Optimizing a squared error against a
zero-filled matrix therefore encodes an assumption that is false by
construction: that absence of interaction means dislike.

Bayesian Personalized Ranking~\cite{rendle2009bpr} replaces that assumption
with a weaker and more defensible one. Rather than predicting an absolute
score, it assumes only that an observed item is preferred to an unobserved
one, and optimizes the probability of that \emph{pairwise} ordering. The
learning signal is a triplet $(u, i, j)$ where $i$ was consumed and $j$ was
not, and the objective rewards $\hat{x}_{ui} > \hat{x}_{uj}$.

\subsection{Scope and claims}

This report accompanies a working repository rather than a novel method. The
claims we make are:

\begin{enumerate}
  \item The pipeline is leakage-free. Every aggregate feature is computed
        point-in-time, and this is enforced by tests that inject a future
        event and assert no feature changes (Section~\ref{sec:features}).
  \item The reported numbers are honest. Hyperparameters for \emph{every}
        model, baselines included, were selected on a validation partition;
        the test partition was read once (Section~\ref{sec:evaluation}).
  \item The served ranking equals the evaluated ranking. Training/serving
        skew is asserted away rather than assumed away
        (Section~\ref{sec:evaluation}).
\end{enumerate}

We explicitly do \emph{not} claim that BPR is the best model for this dataset.
Section~\ref{sec:results} shows it is not.

\subsection{Why the negative result is kept}

Dacrema et al.~\cite{dacrema2019} reproduced eighteen neural recommendation
papers and found that most were beaten by well-tuned non-neural baselines,
largely because the baselines in the original papers had not been tuned at
all. The cheapest way to produce a flattering result for a flagship model is
to compare it against an untuned competitor. We tune the baselines on the same
grid budget as the flagship model, and report the outcome as measured.

\subsection{Repository}

The implementation, infrastructure, and the scripts that generate every table
in this document are available in the accompanying repository. The results
table (Table~\ref{tab:results}) is emitted by
\texttt{scripts/export\_results\_table.py} and the feature catalogue
(Table~\ref{tab:features}) by \texttt{scripts/export\_feature\_tables.py},
so the document cannot drift away from the code.

\section{Data}
\label{sec:data}

\subsection{Source}

The dataset is the Kaggle \emph{Anime Recommendations Database}, containing
approximately $73{,}000$ users, $12{,}000$ titles and $7.8$ million ratings.
Each row is a triple \texttt{(user\_id, anime\_id, rating)} where the rating
is an integer in $[1,10]$, or the sentinel $-1$ meaning the user watched the
title but did not rate it.

The dataset is not redistributed with the repository. A generator produces a
synthetic fixture with the same schema, latent taste clusters and a Pareto
popularity distribution, so that the full pipeline and its tests run without
credentials. All numbers in this report are computed on the fixture; the
qualitative conclusions are what transfer, not the absolute values.

\subsection{From ratings to implicit feedback}

We map explicit ratings to implicit positives:
\begin{equation}
  (u,i) \in \Rset
  \iff
  r_{ui} = -1 \;\;\text{or}\;\; r_{ui} \geq 6 .
  \label{eq:positive}
\end{equation}

Two decisions in Equation~\ref{eq:positive} deserve justification.

\paragraph{The sentinel is a positive.} A user who watched a title without
rating it has revealed consumption, which is exactly the signal BPR models.
Discarding these rows would throw away the majority of the interaction data.

\paragraph{Low ratings are dropped, not treated as negatives.} A rating of $2$
is evidence of \emph{dislike}, which is strictly more information than the
``unobserved'' state BPR assumes for the sampled item $j$. Placing a
known-disliked item into the negative pool would be defensible; placing it
into the positive pool would not. We drop these rows because the sampler
draws negatives from the complement of $\Iset_u^{+}$, and a disliked item
correctly belongs there by default.

\subsection{Core filtering}

Users and items with very few interactions carry almost no collaborative
signal and inflate the catalogue over which ranking metrics are computed. We
apply a $k$-core filter with $k = 5$ on both sides.

The filter is applied \emph{iteratively}. Removing low-activity items reduces
user interaction counts, which can push a previously qualifying user below the
threshold; a single pass would leave the stated invariant violated. The loop
runs until a fixed point, so the output genuinely satisfies
\begin{equation}
  |\Iset_u^{+}| \geq k \;\;\forall u \in \Uset
  \qquad\text{and}\qquad
  |\Uset_i^{+}| \geq k \;\;\forall i \in \Iset .
\end{equation}

\subsection{Temporal ordering and its limitation}

The raw dataset carries \textbf{no timestamps}. This is a real constraint and
we state it plainly rather than inventing one. Point-in-time feature
computation requires an ordering, so we derive a per-user sequence index

\begin{equation}
  t_{ui} = \text{position of } i \text{ within the rows of } u
  \quad\text{(file order)} ,
\end{equation}

and treat it as a weak proxy for chronology.

This proxy is adequate for \emph{enforcing} the leakage discipline --- the
machinery that forbids a feature from reading beyond a cut-off is exercised
and tested regardless of whether the cut-off is a true timestamp. It is
\emph{not} adequate for claiming that the evaluation reproduces a real
production split. On a dataset with genuine timestamps, the same code path
takes the timestamp column instead, and the synthetic fixture carries real
timestamps precisely so that the point-in-time tests exercise that path.

\section{Feature engineering}
\label{sec:features}

Features are declared in \texttt{conf/features.yaml} and executed by a
registry-backed pipeline. The YAML is the single source of truth: a startup
check validates the mapping between declared features and implemented
transforms in \emph{both} directions, so a feature cannot be declared without
an implementation, nor can an implementation silently linger after its
declaration is removed.

% Generated by scripts/export_feature_tables.py -- do not edit by hand.
\begin{table}[t]
  \centering
  \caption{Engineered features. Every aggregate is computed
  point-in-time: only interactions strictly before the prediction
  cut-off contribute to its value.}
  \label{tab:features}
  \begin{tabular}{llll}
    \toprule
    Feature & Family & Transform & Point-in-time \\
    \midrule
    \texttt{user\_activity\_count} & interaction & \texttt{user\_activity\_count} & \checkmark \\
    \texttt{user\_rating\_mean} & interaction & \texttt{user\_rating\_mean} & \checkmark \\
    \texttt{item\_popularity\_rank} & item & \texttt{item\_popularity\_rank} & \checkmark \\
    \texttt{item\_member\_count} & item & \texttt{item\_member\_count} & -- \\
    \texttt{item\_genre\_tokens} & text & \texttt{item\_genre\_tokens} & -- \\
    \texttt{user\_genre\_affinity} & interaction & \texttt{user\_genre\_affinity} & \checkmark \\
    \bottomrule
  \end{tabular}
\end{table}


\subsection{Rationale}

% Generated by scripts/export_feature_tables.py -- do not edit by hand.
\begin{description}
  \item[\texttt{user\_activity\_count}] Separates heavy watchers from casual users. Drives confidence weighting and explains why sparse users need the popularity fallback.

  \item[\texttt{user\_rating\_mean}] Captures user leniency so a score of 7 is read relative to the user's own scale rather than an absolute threshold.

  \item[\texttt{item\_popularity\_rank}] Popularity prior used by the most-popular baseline and as the cold-start fallback ordering for users with no training history.

  \item[\texttt{item\_member\_count}] Catalogue-side audience size. Known before any interaction exists, so it is safe for cold-start items and is not subject to the train-only rule.

  \item[\texttt{item\_genre\_tokens}] Content signal that allows unseen titles to be scored by genre overlap instead of being excluded from recommendation entirely.

  \item[\texttt{user\_genre\_affinity}] Time-decayed taste profile over genres. Recent watches matter more than old ones, so interest fades rather than accumulating forever.

\end{description}


\subsection{Point-in-time computation}

Every aggregate is evaluated against a cut-off $\tau$ using a \emph{strict}
inequality:
\begin{equation}
  \phi(u, \tau) = f\big(\{\, (u,i,t) \in \Rset : t < \tau \,\}\big) .
  \label{eq:pit}
\end{equation}

The strictness matters. With $t \leq \tau$, the interaction being predicted
contributes to its own feature value, and the model receives a trace of the
label. This is the most common leakage bug in recommender pipelines because
it produces excellent offline metrics and no error message.

Concretely, user activity count and the rating mean are
\begin{align}
  \mathrm{activity}(u,\tau) &= \big|\{\, (u,i,t) \in \Rset : t < \tau \,\}\big| ,\\
  \overline{r}(u,\tau)      &= \frac{1}{\mathrm{activity}(u,\tau)}
                               \sum_{(u,i,t) \in \Rset,\, t < \tau} r_{ui} ,
\end{align}
with $\overline{r}(u,\tau)$ defined as a configured fill value when the
denominator is zero.

Genre affinity is a time-decayed profile over the genre vocabulary $G$. For
genre $g$,
\begin{equation}
  \mathrm{aff}(u, g, \tau)
  = \frac{\sum_{(u,i,t) \in \Rset,\, t<\tau} \lambda^{\,\tau - t}\, \mathbb{1}[g \in G_i]}
         {\sum_{(u,i,t) \in \Rset,\, t<\tau} \lambda^{\,\tau - t}} ,
  \qquad \lambda \in (0, 1] ,
  \label{eq:affinity}
\end{equation}
so that recent consumption dominates a user's represented taste.

Item popularity is a rank, not a raw count:
\begin{equation}
  \mathrm{poprank}(i, \tau) = \big|\{\, i' \in \Iset : c_{i'}(\tau) > c_{i}(\tau) \,\}\big| + 1 ,
  \qquad
  c_i(\tau) = \big|\{\,(u,i,t) \in \Rset : t < \tau\,\}\big| .
\end{equation}
Ranking is deliberate: raw counts follow a Pareto distribution whose scale
shifts as the catalogue grows, whereas the rank is stable under that growth.

\subsection{Fitting discipline}

Statistics that require fitting --- popularity tables, encodings, normalizer
parameters --- are fitted on the \emph{training partition only} and then
applied to validation and test. Fitting on the full dataset is leakage even
when the per-row computation respects the cut-off, because the fitted
parameters summarise the held-out rows.

\subsection{Leakage controls}

Three classes of test enforce the above, in
\texttt{tests/test\_features\_leakage.py}:

\begin{enumerate}
  \item \textbf{Future-event injection.} An interaction is appended
        \emph{after} the cut-off and every feature is recomputed. Any change
        in any value fails the test. This is the direct executable form of
        Equation~\ref{eq:pit}.
  \item \textbf{Fit-partition isolation.} Fitted tables computed on the
        training partition are asserted identical whether or not the
        validation and test rows are present in the input.
  \item \textbf{Label destruction.} Item assignments in the training
        partition are randomly permuted, destroying the user--item
        association while preserving every marginal distribution. A model
        retrained on the corrupted data must collapse towards the random
        baseline. If it does not, some feature is carrying the answer.
\end{enumerate}

The third test is the one that catches leakage the structural checks miss,
because it makes no assumption about \emph{which} feature is responsible.

\subsection{A dtype bug found by these tests}

The future-event injection test initially failed with an index dtype mismatch
(\texttt{int32} versus \texttt{int64}) rather than a value difference. The
cause was that \texttt{pandas} promotes index dtypes differently depending on
how a frame was assembled, so two logically identical feature frames compared
unequal. The fix normalises every user-keyed series to an \texttt{int64}
index at construction. We record it here because it is a genuine robustness
defect that would have surfaced as a confusing production failure, not a
test artifact that was patched away.

\section{Methodology}
\label{sec:methodology}

\subsection{The BPR criterion}

BPR starts from a per-user total order $>_u$ over items and assumes only that
observed items are preferred to unobserved ones. The training set of triplets
is
\begin{equation}
  D_{\Rset} = \{\, (u, i, j) \mid i \in \Iset_u^{+},\; j \in \Iset \setminus \Iset_u^{+} \,\} .
\end{equation}

Assuming independence across users and triplets, the posterior over
parameters given the orderings is
\begin{equation}
  p(\Theta \mid >_u) \propto p(>_u \mid \Theta)\, p(\Theta) ,
\end{equation}
and with the pairwise likelihood modelled as
$p(i >_u j \mid \Theta) = \sigma(\hat{x}_{uij})$ and a Gaussian prior
$\Theta \sim \mathcal{N}(0, \lambda^{-1} I)$, maximizing the posterior is
equivalent to maximizing

\begin{equation}
  \mathrm{BPR\text{-}Opt}
  = \sum_{(u,i,j) \in D_{\Rset}} \ln \sigma(\hat{x}_{uij})
    \;-\; \lambda \lVert \Theta \rVert^{2} ,
  \label{eq:bpropt}
\end{equation}

where the pairwise score difference is
\begin{equation}
  \hat{x}_{uij} = \hat{x}_{ui} - \hat{x}_{uj} .
\end{equation}

\subsection{Matrix factorization scorer}

We instantiate the scorer as biased matrix factorization:
\begin{equation}
  \hat{x}_{ui} = \pvec_u^{\top} \qvec_i + b_i
  = \sum_{f=1}^{d} p_{uf} q_{if} + b_i .
  \label{eq:mf}
\end{equation}

The user bias is deliberately absent. A per-user constant shifts every item
score for that user by the same amount and therefore cannot change the
induced ranking; including it would only add an unidentifiable parameter.
The item bias survives because it \emph{does} reorder items, acting as a
learned popularity prior. Note also that $b_i$ cancels from
$\hat{x}_{uij}$ only when $i = j$, which never occurs.

\subsection{Gradients}

Write $\delta = \sigma(-\hat{x}_{uij})$, which follows from
$\frac{\partial}{\partial z}\ln\sigma(z) = \sigma(-z)$. For a single triplet
the gradients of Equation~\ref{eq:bpropt} are

\begin{align}
  \frac{\partial}{\partial \pvec_u} &= \delta\,(\qvec_i - \qvec_j) - \lambda_u \pvec_u ,\\
  \frac{\partial}{\partial \qvec_i} &= \delta\,\pvec_u - \lambda_i \qvec_i ,\\
  \frac{\partial}{\partial \qvec_j} &= -\delta\,\pvec_u - \lambda_i \qvec_j ,\\
  \frac{\partial}{\partial b_i}     &= \delta - \lambda_b b_i ,\\
  \frac{\partial}{\partial b_j}     &= -\delta - \lambda_b b_j .
  \label{eq:grads}
\end{align}

The factor $\delta$ has a useful reading: it is large when the model
currently ranks the triplet \emph{wrongly} and tends to zero once the pair is
confidently ordered, so converged triplets stop contributing gradient.

Parameters are updated by ascent with step size $\eta$ decayed
geometrically per epoch, $\eta_{e} = \eta_0 \gamma^{e}$.

\subsection{Vectorization and why it changed the conclusions}

A literal reading of Equation~\ref{eq:grads} gives one update per triplet.
In NumPy that costs roughly $70$ seconds for $40$ epochs on the fixture,
which makes a hyperparameter grid search impractical --- and an untuned
flagship model compared against baselines would have reproduced exactly the
methodological failure criticised in Section~\ref{sec:introduction}.

We therefore apply the same gradients to mini-batches of $B$ triplets. One
subtlety makes this non-trivial: a user or item may appear several times
within a batch, and NumPy's fancy-indexed assignment
\texttt{A[idx] += v} keeps only the \emph{last} write for duplicated indices,
silently discarding gradients. We use \texttt{np.add.at}, which accumulates.
The resulting implementation runs the same $40$ epochs in $2.3$ seconds, a
$\sim 30\times$ speed-up, which is what made the grid search in
Section~\ref{sec:results} affordable.

\subsection{Negative sampling}

For each positive $(u,i)$ the sampler draws $j$ uniformly from
$\Iset \setminus \Iset_u^{+}$. Rejection sampling is the natural approach,
but testing membership against a per-user Python set inside the training loop
reintroduces the interpreter overhead that vectorization removed, and a dense
$m \times n$ boolean mask would need roughly $876$\,MB on the full dataset.

We instead encode each observed pair as a single integer key
\begin{equation}
  \kappa(u,i) = u \cdot n + i ,
\end{equation}
store the sorted key array, and test candidate membership with a vectorized
binary search. Memory is $O(|\Rset|)$ rather than $O(mn)$, and collisions are
resampled for a bounded number of rounds. A test asserts zero false negatives
across $340{,}000$ draws.

\subsection{Baselines}

\paragraph{Most-popular.} Ranks all users by the same global interaction
count. It is non-personalised by construction and serves as the floor that
any personalised model must clear.

\paragraph{Item-$k$NN.} Cosine similarity between binary item vectors,
\begin{equation}
  s_{ij} = \frac{\lvert \Uset_i^{+} \cap \Uset_j^{+} \rvert}
                {\sqrt{\lvert \Uset_i^{+} \rvert}\sqrt{\lvert \Uset_j^{+} \rvert}} ,
\end{equation}
truncated to the $k$ nearest neighbours per item, with user score
\begin{equation}
  \hat{x}_{ui} = \sum_{j \in \Iset_u^{+}} s_{ij} .
\end{equation}
This is the method most Kaggle notebooks for this dataset actually use, so it
is the bar that justifies the project.

\paragraph{Random.} A seeded permutation per user. It calibrates the metric
scale and, with full catalogue coverage, demonstrates that coverage alone is
not a quality signal.

\section{Evaluation protocol}
\label{sec:evaluation}

\subsection{Partitioning}

We use leave-last-one-out per user. Ordering each user's interactions by the
temporal proxy of Section~\ref{sec:data}, the final interaction forms the test
instance, the penultimate one the validation instance, and everything earlier
is training data.

Users with fewer than $k_{\min} + 2$ interactions cannot yield both held-out
instances without starving their own training history; these users contribute
entirely to the training partition and are not evaluated. Reporting metrics
over users whose training history is empty would measure the popularity prior
rather than personalisation.

\subsection{Candidate set}

For each evaluated user we rank the \emph{entire} catalogue minus that user's
training history:
\begin{equation}
  C_u = \Iset \setminus \Iset_u^{+,\text{train}} .
\end{equation}

Excluding the history is not a detail. Without it, a model is rewarded for
re-recommending already-consumed titles, which inflates every metric and
corresponds to a product behaviour nobody wants. A test constructs a model
that returns only the user's history and asserts it scores exactly zero.

We rank the full catalogue rather than sampling $100$ negatives per user.
Sampled-negative evaluation is common but is known to produce rankings of
methods that disagree with the full-catalogue result~\cite{krichene2020}.

\subsection{Metrics}

With exactly one relevant item $i^{*}_u$ per evaluated user and
$\mathrm{rank}_u \in \{1, \dots, |C_u|\}$ its position in the predicted list,
the cut-off metrics over the evaluated user set $\Uset_{\text{eval}}$ are:

\begin{align}
  \mathrm{NDCG@}K &= \frac{1}{|\Uset_{\text{eval}}|} \sum_{u}
    \frac{\mathbb{1}[\mathrm{rank}_u \leq K]}{\log_2(\mathrm{rank}_u + 1)} ,
    \label{eq:ndcg} \\[0.3em]
  \mathrm{Recall@}K &= \frac{1}{|\Uset_{\text{eval}}|} \sum_{u}
    \mathbb{1}[\mathrm{rank}_u \leq K] , \\[0.3em]
  \mathrm{MRR@}K &= \frac{1}{|\Uset_{\text{eval}}|} \sum_{u}
    \frac{\mathbb{1}[\mathrm{rank}_u \leq K]}{\mathrm{rank}_u} , \\[0.3em]
  \mathrm{Coverage@}K &= \frac{1}{|\Iset|}
    \Big\lvert \bigcup_{u \in \Uset_{\text{eval}}} \mathrm{top}_K(u) \Big\rvert .
\end{align}

Equation~\ref{eq:ndcg} uses the general NDCG definition specialised to a
single relevant item, where the ideal DCG equals $1$ and the normalisation
therefore disappears. Under this protocol Recall@$K$ coincides numerically
with hit rate; both are reported so the table stays comparable with
multi-positive projects.

Coverage is reported because it is the business reason to leave a popularity
list, but it must be read alongside accuracy: the random baseline attains
coverage $1.0$ while being useless.

\subsection{Model selection}

All hyperparameters --- for the baselines as well as for BPR --- are selected
by grid search on the \emph{validation} partition. The test partition is read
once, to produce Table~\ref{tab:results}.

This is stated explicitly because the alternative is both common and
invisible in a results table. Selecting on test turns the reported number
into an optimistically biased estimate whose gap from reality grows with the
size of the grid.

\subsection{Training/serving parity}

Offline metrics are worthless if production serves a different ranking. The
serving path precomputes top-$N$ lists into DynamoDB, and
\texttt{tests/test\_features\_parity.py} asserts that the stored list for each
user is element-wise identical to the list produced by the same fitted model
during evaluation, that no served list intersects the user's history, and
that the payload is JSON-serialisable --- NumPy scalars are a frequent and
late-surfacing cause of rejected writes.

\section{Results}
\label{sec:results}

% Generated by scripts/export_results_table.py -- do not edit by hand.
\begin{table}[t]
  \centering
  \caption{Top-$K$ ranking quality on the held-out test partition ($K=10$).
  Hyperparameters for every model were selected on the validation
  partition; the test partition was read once.}
  \label{tab:results}
  \begin{tabular}{lrrrr}
    \toprule
    Model & NDCG@10 & Recall@10 & MRR@10 & Coverage \\
    \midrule
    ItemKNN & \textbf{0.1305} & 0.2483 & 0.0952 & 0.4969 \\
    BPR & 0.1248 & 0.2339 & 0.0922 & 0.3937 \\
    Popularity & 0.0761 & 0.1448 & 0.0556 & 0.0750 \\
    Random & 0.0150 & 0.0334 & 0.0095 & 1.0000 \\
    \bottomrule
  \end{tabular}
\end{table}


\subsection{Reading the table}

Both personalised models clear the most-popular baseline by a wide margin,
and all three clear random by roughly an order of magnitude, so the pipeline
is learning genuine collaborative structure.

The flagship model, however, loses. Tuned item-$k$NN reaches NDCG@10 of
$0.1305$ against $0.1248$ for tuned BPR, a relative gap of about $4.6\%$ in
favour of the baseline. Recall@10 and MRR@10 order the two models the same
way, so the result is not an artifact of a single metric's discount function.

\subsection{Hyperparameter search}

The BPR grid covered $d \in \{16, 32, 64, 128\}$,
$\eta_0 \in \{0.02, 0.05, 0.1, 0.2\}$,
$\lambda \in \{0.005, 0.02, 0.05, 0.1, 0.2\}$ and epoch budgets from $20$ to
$200$, selected on validation NDCG@10. The optimum sits at $d = 32$,
$\eta_0 = 0.2$, $\lambda = 0.05$.

Two features of the search are worth recording. First, validation performance
is \emph{non-monotonic} in the epoch budget: it peaks and then declines while
the training objective keeps improving, which is overfitting of the pairwise
criterion rather than of a reconstruction error. Second, performance is far
more sensitive to $\lambda$ than to $d$ --- the swing across the
regularization axis is roughly $0.05$ NDCG, against roughly $0.005$ across
the factor axis. The item-$k$NN neighbourhood size was swept over
$k \in \{10, 20, 50, 100, 200, |\Iset|\}$ and is remarkably flat, varying by
under $0.004$ NDCG across the entire range.

That flatness is itself informative: the baseline is near-optimal almost
regardless of its one hyperparameter, while BPR requires a search to become
competitive at all.

\subsection{Why the baseline wins here}

The result is consistent with the data regime rather than with a defect in
the implementation, and three observations support that reading.

\paragraph{Scale.} Matrix factorization amortises a shared low-rank structure
across many observations. The evaluated fixture has $\sim 900$ users and
$320$ items. Repeating the comparison at roughly five times the scale
($\sim 5{,}000$ users, $1{,}500$ items, $95{,}000$ training interactions) did
not reverse the ordering, but it did compress the gap, which is the direction
the scale argument predicts.

\paragraph{Density.} The $k$-core filter leaves a comparatively dense
interaction matrix. Co-occurrence counts are therefore well estimated
directly, and the cosine similarities that item-$k$NN consumes are low
variance. Neighbourhood methods are strongest exactly in this regime.

\paragraph{Intrinsic rank.} The latent structure is a small number of taste
clusters crossed with a popularity prior. Item-$k$NN represents
within-cluster co-consumption essentially exactly, whereas BPR must recover
the same structure through a rank-$d$ approximation fitted by sampled SGD.
There is little headroom above a model that already captures the generating
process.

\subsection{Coverage}

BPR covers $39\%$ of the catalogue against $50\%$ for item-$k$NN and $7.5\%$
for most-popular. The popularity baseline's figure quantifies the
concentration problem that motivates personalisation: it is structurally
incapable of surfacing more than a handful of titles. Random attains $100\%$
coverage with near-zero accuracy, which is the reminder that coverage is a
constraint to satisfy, not an objective to maximise.

\subsection{Threats to validity}

\begin{enumerate}
  \item \textbf{Synthetic fixture.} The reported numbers come from generated
        data. The generating process was not chosen to favour either model,
        but it is simpler than reality, and simple structure favours the
        method with the stronger inductive bias.
  \item \textbf{Proxy ordering.} The split uses file order as a chronology
        proxy (Section~\ref{sec:data}). A true temporal split would be
        harder, and would likely lower all reported numbers.
  \item \textbf{Single seed.} Metrics are reported for one seed. The gap
        between BPR and item-$k$NN is small enough that a confidence interval
        over seeds would be needed before calling it decisive.
  \item \textbf{Offline only.} NDCG on held-out interactions is a proxy for
        user satisfaction and cannot capture novelty, diversity or
        presentation effects.
\end{enumerate}

\section{Conclusion}
\label{sec:conclusion}

We implemented BPR matrix factorization end to end: implicit-feedback
construction, iterative core filtering, point-in-time feature engineering,
vectorized triplet SGD, full-catalogue ranking evaluation, and a
parity-tested serving path deployed by Terraform against a local AWS
emulator.

The measured outcome is that a tuned item-$k$NN baseline edges out tuned BPR
on this data (NDCG@10 $0.1305$ versus $0.1248$). We report it because the
alternative --- tuning until the flagship model wins --- is the specific
failure mode that Dacrema et al.~\cite{dacrema2019} documented across a
large body of published recommender work, and avoiding it was an explicit
goal of this project.

\subsection{What transfers}

The reusable artifacts of this project are not the model weights:

\begin{itemize}
  \item \textbf{A leakage harness.} Future-event injection, fit-partition
        isolation, and label destruction, as executable tests rather than
        prose. The injection test surfaced a genuine dtype-promotion defect
        in the feature layer.
  \item \textbf{A parity harness.} Asserting that the served ranking equals
        the evaluated ranking closes the gap where good offline numbers
        quietly fail to reach production.
  \item \textbf{An honest selection protocol.} Validation-only tuning, with
        the same grid budget extended to the baselines.
\end{itemize}

\subsection{Future work}

\paragraph{Popularity-weighted sampling.} Drawing negatives
$j$ with probability proportional to $c_j^{\alpha}$ rather than uniformly
produces harder triplets as training progresses, and is the most likely
single change to close the observed gap.

\paragraph{Confidence intervals.} Repeating the comparison across seeds and
reporting intervals would establish whether the $0.006$ NDCG gap is
significant.

\paragraph{Hybrid scoring.} The content features in
Table~\ref{tab:features} are currently consumed only for cold-start fallback.
Folding them into the scorer, $\hat{x}_{ui} = \pvec_u^{\top}\qvec_i + b_i +
\mathbf{w}^{\top}\boldsymbol{\phi}_{ui}$, would let the model exploit genre
affinity directly.

\paragraph{Real timestamps.} Evaluating on a dataset with genuine event times
would replace the chronology proxy and make the temporal split defensible as
a simulation of production.


\bibliographystyle{plain}
\bibliography{references}

\end{document}
